\subsection{Chapter Summary}

This chapter presented the theoretical background required for the study of OTFS and OTFS-IM systems. The fundamental characteristics of wireless communication channels were first introduced, including channel impairments, channel classification, multipath propagation, delay spread, Doppler effects, and doubly selective fading channels.

The chapter then reviewed the principles of OFDM, including its transceiver structure, advantages, and limitations in high-mobility environments. Although OFDM provides high spectral efficiency and simple frequency-domain equalization, its performance degrades significantly when severe Doppler spread destroys subcarrier orthogonality and introduces inter-carrier interference.

The delay-Doppler domain representation was also discussed as an alternative way to characterize time-varying wireless channels. This representation describes wireless propagation in terms of delay and Doppler shifts, providing a more suitable foundation for understanding OTFS modulation. Finally, the basic concept of Index Modulation was introduced, including its principle, OFDM-IM structure, advantages, and challenges.

Overall, this chapter established the necessary theoretical foundation for the following chapters. Based on these concepts, Chapter 3 will present the principles and mathematical formulation of OTFS modulation in more detail.

\newpage